% GENERATED FILE. Edit canonical lessons, then run npm run generate:books.
% canonical-source-hash: fnv1a64:79ee202f9c8aae19
% canonical-lessons: MR-C552-phiskatne, MR-C552-bargalne, MR-C552-dhasalne, MR-C552-kolmadne, MR-C552-sutne

\chapter{When Plans Fail}
\label{ch:mr-when-plans-fail}
\hlchaptermodality{552}

\begin{chapteropening}
\textbf{By the end of this chapter:} \emph{I can say how plans fall apart and things come loose.}

Complete the last lesson of chapter 552: I can say how plans fall apart and things come loose.
\end{chapteropening}

\section[\texorpdfstring{\mr{फिसकटणे} (phiskaṭṇe)}{फिसकटणे (phiskaṭṇe)}]{\mr{फिसकटणे} (phiskaṭṇe) — to fall apart; to fall through}

\label{lesson:MR-C552-phiskatne}



\begin{quote}
Before the new one: say the Marathi for reserved, then the Marathi for unforeseen.
\end{quote}

\subsection*{You'll want to know: \mr{फिसकटणे}}
\textbf{\mr{फिसकटणे}} — \emph{phiskaṭṇe} — \textquotedblleft{}to fall apart; to fall through\textquotedblright{}.

\begin{quote}
\mr{पावसामुळे} \mr{आमचा} \mr{सहलीचा} \mr{बेत} \mr{फिसकटला}.
\end{quote}

\begin{quote}
\emph{pāvsāmuḷe āmcā sahlīcā bet phiskaṭlā.}
\end{quote}

\begin{quote}
Because of the rain our plan for a trip fell through.
\end{quote}

Past (he …): \mr{तो} \mr{सौदा} \mr{फिसकटला}, \emph{to saudā phiskaṭlā}.

\begin{cousinweb}
\textbf{\mr{फिसकटणे}} (\emph{phiskaṭṇe}) is inherited: CDIAL lists Marathi \emph{phiskaṭṇẽ}, \textquotedblleft{}to fall away\textquotedblright{}, under \emph{*phiss}, \textquotedblleft{}slip\textquotedblright{} (CDIAL 9080).
\end{cousinweb}

\begin{grammarlens}[title={What you've built}]
One more word for saying what happens if.
\end{grammarlens}

\subsection*{Your turn}
\begin{itemize}
\raggedright
  \item \emph{Say it:} \emph{phiskaṭṇe}
  \item \emph{Say it:} \emph{phiskaṭṇe}, once more
  \item \emph{Say it:} \emph{pāvsāmuḷe āmcā sahlīcā bet phiskaṭlā.} — Because of the rain our plan for a trip fell through.
  \item \emph{Recall:} say the Marathi for reserved, then the Marathi for unforeseen, then say \emph{phiskaṭṇe} again
\end{itemize}

\begin{tcolorbox}[breakable,colback=teal!4,colframe=teal!35!black,title={Before you move on}]
What does \mr{फिसकटणे} mean? (\textquotedblleft{}To fall apart; to fall through\textquotedblright{}.) Say it once more.
\end{tcolorbox}
\section[\texorpdfstring{\mr{बारगळणे} (bārgaḷṇe)}{बारगळणे (bārgaḷṇe)}]{\mr{बारगळणे} (bārgaḷṇe) — to fizzle out}

\label{lesson:MR-C552-bargalne}



\begin{quote}
Before the new one: say the Marathi for unforeseen, then the Marathi for to fall apart.
\end{quote}

\subsection*{You'll want to know: \mr{बारगळणे}}
\textbf{\mr{बारगळणे}} — \emph{bārgaḷṇe} — \textquotedblleft{}to fizzle out\textquotedblright{}.

\begin{quote}
\mr{पैसे} \mr{नसल्यामुळे} \mr{नवीन} \mr{दुकानाची} \mr{योजना} \mr{बारगळली}.
\end{quote}

\begin{quote}
\emph{paise naslyāmuḷe navīn dukānācī yojnā bārgaḷlī.}
\end{quote}

\begin{quote}
For lack of money the plan for a new shop fizzled out.
\end{quote}

Past (he …): \mr{तो} \mr{प्रयत्न} \mr{बारगळला}, \emph{to prayatna bārgaḷlā}.

\begin{cousinweb}
\textbf{\mr{बारगळणे}} (\emph{bārgaḷṇe}) is a Marathi verb for a plan that comes to nothing; its origin is not settled.
\end{cousinweb}

\begin{grammarlens}[title={What you've built}]
One more word for saying what happens if.
\end{grammarlens}

\subsection*{Your turn}
\begin{itemize}
\raggedright
  \item \emph{Say it:} \emph{bārgaḷṇe}
  \item \emph{Say it:} \emph{bārgaḷṇe}, once more
  \item \emph{Say it:} \emph{paise naslyāmuḷe navīn dukānācī yojnā bārgaḷlī.} — For lack of money the plan for a new shop fizzled out.
  \item \emph{Recall:} say the Marathi for unforeseen, then the Marathi for to fall apart, then say \emph{bārgaḷṇe} again
\end{itemize}

\begin{tcolorbox}[breakable,colback=teal!4,colframe=teal!35!black,title={Before you move on}]
What does \mr{बारगळणे} mean? (\textquotedblleft{}To fizzle out\textquotedblright{}.) Say it once more.
\end{tcolorbox}
\section[\texorpdfstring{\mr{ढासळणे} (ḍhāsaḷṇe)}{ढासळणे (ḍhāsaḷṇe)}]{\mr{ढासळणे} (ḍhāsaḷṇe) — to crumble; to cave in}

\label{lesson:MR-C552-dhasalne}



\begin{quote}
Before the new one: say the Marathi for to fall apart, then the Marathi for to fizzle out.
\end{quote}

\subsection*{You'll want to know: \mr{ढासळणे}}
\textbf{\mr{ढासळणे}} — \emph{ḍhāsaḷṇe} — \textquotedblleft{}to crumble; to cave in\textquotedblright{}.

\begin{quote}
\mr{जोरदार} \mr{पावसात} \mr{जुन्या} \mr{किल्ल्याची} \mr{भिंत} \mr{ढासळली}.
\end{quote}

\begin{quote}
\emph{jordār pāvsāt junyā killyācī bhint ḍhāsaḷlī.}
\end{quote}

\begin{quote}
In the heavy rain the old fort's wall crumbled.
\end{quote}

Past (he …): \mr{तो} \mr{किल्ला} \mr{ढासळला}, \emph{to killā ḍhāsaḷlā}.

\begin{cousinweb}
\textbf{\mr{ढासळणे}} (\emph{ḍhāsaḷṇe}) is a tadbhava: CDIAL lists Marathi \emph{dhāsaḷṇẽ}, \textquotedblleft{}to fall\textquotedblright{}, under Sanskrit \emph{dhvaṁsati}, \textquotedblleft{}falls to pieces\textquotedblright{} (CDIAL 6896).
\end{cousinweb}

\begin{grammarlens}[title={What you've built}]
One more word for saying what happens if.
\end{grammarlens}

\subsection*{Your turn}
\begin{itemize}
\raggedright
  \item \emph{Say it:} \emph{ḍhāsaḷṇe}
  \item \emph{Say it:} \emph{ḍhāsaḷṇe}, once more
  \item \emph{Say it:} \emph{jordār pāvsāt junyā killyācī bhint ḍhāsaḷlī.} — In the heavy rain the old fort's wall crumbled.
  \item \emph{Recall:} say the Marathi for to fall apart, then the Marathi for to fizzle out, then say \emph{ḍhāsaḷṇe} again
\end{itemize}

\begin{tcolorbox}[breakable,colback=teal!4,colframe=teal!35!black,title={Before you move on}]
What does \mr{ढासळणे} mean? (\textquotedblleft{}To crumble; to cave in\textquotedblright{}.) Say it once more.
\end{tcolorbox}
\section[\texorpdfstring{\mr{कोलमडणे} (kolmaḍṇe)}{कोलमडणे (kolmaḍṇe)}]{\mr{कोलमडणे} (kolmaḍṇe) — to break down; to collapse}

\label{lesson:MR-C552-kolmadne}



\begin{quote}
Before the new one: say the Marathi for to fizzle out, then the Marathi for to crumble.
\end{quote}

\subsection*{You'll want to know: \mr{कोलमडणे}}
\textbf{\mr{कोलमडणे}} — \emph{kolmaḍṇe} — \textquotedblleft{}to break down; to collapse\textquotedblright{}.

\begin{quote}
\mr{वडील} \mr{आजारी} \mr{पडले} \mr{आणि} \mr{घराचं} \mr{सगळं} \mr{गणित} \mr{कोलमडलं}.
\end{quote}

\begin{quote}
\emph{vaḍīl ājārī paḍle āṇi gharācaṁ sagḷaṁ gaṇit kolmaḍlaṁ.}
\end{quote}

\begin{quote}
Father fell ill and the family's whole budget collapsed.
\end{quote}

Past (he …): \mr{तो} \mr{कोलमडला}, \emph{to kolmaḍlā}.

\begin{cousinweb}
\textbf{\mr{कोलमडणे}} (\emph{kolmaḍṇe}) is a Marathi verb for a person or a plan that gives way; its origin is not settled.
\end{cousinweb}

\begin{grammarlens}[title={What you've built}]
One more word for saying what happens if.
\end{grammarlens}

\subsection*{Your turn}
\begin{itemize}
\raggedright
  \item \emph{Say it:} \emph{kolmaḍṇe}
  \item \emph{Say it:} \emph{kolmaḍṇe}, once more
  \item \emph{Say it:} \emph{vaḍīl ājārī paḍle āṇi gharācaṁ sagḷaṁ gaṇit kolmaḍlaṁ.} — Father fell ill and the family's whole budget collapsed.
  \item \emph{Recall:} say the Marathi for to fizzle out, then the Marathi for to crumble, then say \emph{kolmaḍṇe} again
\end{itemize}

\begin{tcolorbox}[breakable,colback=teal!4,colframe=teal!35!black,title={Before you move on}]
What does \mr{कोलमडणे} mean? (\textquotedblleft{}To break down; to collapse\textquotedblright{}.) Say it once more.
\end{tcolorbox}
\section[\texorpdfstring{\mr{सुटणे} (suṭṇe)}{सुटणे (suṭṇe)}]{\mr{सुटणे} (suṭṇe) — to come loose; to be let off}

\label{lesson:MR-C552-sutne}



\begin{quote}
Before the new one: say the Marathi for to crumble, then the Marathi for to break down.
\end{quote}

\subsection*{You'll want to know: \mr{सुटणे}}
\textbf{\mr{सुटणे}} — \emph{suṭṇe} — \textquotedblleft{}to come loose; to be let off\textquotedblright{}.

\begin{quote}
\mr{गाठ} \mr{नीट} \mr{बांधली} \mr{नव्हती}, \mr{म्हणून} \mr{दोरी} \mr{सुटली}.
\end{quote}

\begin{quote}
\emph{gāṭh nīṭ bāndhlī navhatī, mhaṇūn dorī suṭlī.}
\end{quote}

\begin{quote}
The knot was not tied well, so the rope came loose.
\end{quote}

Past (he …): \mr{तो} \mr{सुटला}, \emph{to suṭlā}.

Also \mr{शाळा} \mr{सुटली}, school let out, and \mr{प्रश्न} \mr{सुटला}, the problem was solved.

\begin{cousinweb}
\textbf{\mr{सुटणे}} (\emph{suṭṇe}) is inherited: CDIAL lists Marathi \emph{suṭā}, \textquotedblleft{}loosened\textquotedblright{}, under \emph{*kṣuṭyatē}, \textquotedblleft{}is released\textquotedblright{} (CDIAL 3707).
\end{cousinweb}

\begin{grammarlens}[title={What you've built}]
One more word for saying what happens if.
\end{grammarlens}

\subsection*{Your turn}
\begin{itemize}
\raggedright
  \item \emph{Say it:} \emph{suṭṇe}
  \item \emph{Say it:} \emph{suṭṇe}, once more
  \item \emph{Say it:} \emph{gāṭh nīṭ bāndhlī navhatī, mhaṇūn dorī suṭlī.} — The knot was not tied well, so the rope came loose.
  \item \emph{Recall:} say the Marathi for to crumble, then the Marathi for to break down, then say \emph{suṭṇe} again
\end{itemize}

\begin{tcolorbox}[breakable,colback=teal!4,colframe=teal!35!black,title={Before you move on}]
What does \mr{सुटणे} mean? (\textquotedblleft{}To come loose; to be let off\textquotedblright{}.) Say it once more.
\end{tcolorbox}
